\section{Operationalization of the Concentration-Gap~Conjecture}
\label{sec:conjecture_op}

Theorem~\ref{thm:deployment_safety} depends on
\citep[Microfoundation]{lasser2026micro}'s Concentration-Gap~Conjecture
(optimization pressure correlates with gap exploitation) through invariant
$I_5$ (HHI ceiling).  This section summarizes the dependency
structure, points the reader at the companion
paper~\citep{lasser2026foundations} for the formal discharge, and
discusses what remains conditional after that discharge.

\subsection{The dependency, before the companion paper}

\citep[Microfoundation]{lasser2026micro}'s Concentration-Gap~Conjecture
states informally: optimization pressure on the proxy $\Pproxy =
\volL$ correlates with exploitation of the proxy-truth gap
$\epsnonres$.  The conjecture's deployment relevance: when the
system is in the optimization-pressure regime
$\mathcal{R}_{\mathrm{press}}$, predicted gap exploitation
invalidates the static Goodhart bound's intensivity property.
Theorem~\ref{thm:deployment_safety}'s Layer~1 binding requires the
deployment to lie outside $\mathcal{R}_{\mathrm{press}}$.

The HHI ceiling $I_5$ provides an operational characterization of
``outside the pressure regime'' through trade-flow concentration.
Lemma~\ref{lem:3_forward} establishes the forward implication
($\HHI > H^*$ implies the trade-flow concentration prerequisites
of $\mathcal{R}_{\mathrm{press}}$); the converse-direction
sufficiency required by the theorem ($\HHI < H^*$ ruling out the
pressure regime) is supplied by the scoped Concentration-Gap
selection theorem of \citep[Theorem~2]{lasser2026foundations}
under Assumption~\ref{ass:concgap_conditions}'s structural
conditions.

\subsection{The companion-paper discharge}

The companion paper~\citep{lasser2026foundations} discharges a
\emph{scoped} version of the Concentration-Gap~Conjecture under
four concrete structural conditions:

\begin{itemize}
\item \textbf{(REP) Representativeness}: the counterparty
  population's mean utility is bounded close to the
  welfare-relevant truth $W$.
\item \textbf{(DISP) Bounded dispersion}: counterparty utility
  deviations have bounded population variance.
\item \textbf{(COAL) Coalition closure}: counterparty correlations
  above an audit threshold are partitioned into coalitions, with
  the post-partition counterparty cardinality $N$ deployment-class
  bounded and the latent-coalition residual bounded operationally.
\item \textbf{(EMB) Lipschitz embedding compatibility}: the
  embedding $\phi: \mathcal{V}_\mathrm{utility} \to \mathcal{V}$
  is affine isometric (or bi-Lipschitz with documented
  constants) on the operationally active subspace.
\end{itemize}

\citep[Theorem~2]{lasser2026foundations} establishes that under
(REP) + (DISP) + (COAL) and Lipschitz embedding compatibility
(EMB), the proxy-truth Goodhart slack is bounded by an explicit
$\chi^2$-divergence quantity that controls $\HHI$ via $\chi^2(w
\,\|\, \mu) = N \cdot \HHI(w) - 1$ for uniform base measure $\mu$
on a (COAL)-bounded post-partition counterparty population.
Layer~1 of Theorem~\ref{thm:deployment_safety} binds via this
structural result.

The four structural conditions are operationally auditable;
audit procedures are specified in
\citep[\S 6]{lasser2026foundations} and integrate with this
paper's \S\ref{sec:deployment_tooling} via the audit hooks
listed there.  Assumption~\ref{ass:concgap_conditions} states the
conjunction (REP) + (DISP) + (COAL) + (EMB) as the load-bearing
structural content of the deployment claim's Layer~1 binding.

\subsection{What remains conditional}

The companion paper discharges a scoped version of the
Concentration-Gap~Conjecture under named structural conditions.
What remains:

\begin{itemize}
\item \textbf{The universal model-free Concentration-Gap~Conjecture
  is not proved.}  \citep[\S 7]{lasser2026foundations} argues that
  this conjecture is not provable from current GFM machinery
  without adding a formal optimizer/selection model.  The scoped
  discharge is sufficient for the deployment claim's purposes;
  the universal form is a research-program target.

\item \textbf{The $g$-level reverse direction is not established.}
  \citep[Theorem~2]{lasser2026foundations} bounds the proxy-truth
  norm $\|\Pproxy - \Ttruth\|$ from below under high HHI plus
  separation conditions, but does not transfer this to a lower
  bound on $|g(\Ttruth) - g(\Pproxy)|$ without additional
  inverse-Lipschitz structure on $g$ (constant $g$
  counterexample).  For deployment-claim purposes, the norm-level
  reverse is sufficient: the detection layer (Layer~2) operates
  on the four-channel observable structure, not directly on
  alignment-property values.

\item \textbf{The latent-coalition residual} $\eta_{\mathrm{latent}}$
  is bounded operationally rather than structurally.  Deployments
  must calibrate audit detection-power to bound
  $\eta_{\mathrm{latent}}$; this is operational tooling
  (\S\ref{sec:deployment_tooling}) rather than internal to the
  algebraic theorem.
\end{itemize}

These open items do not threaten the deployment claim under the
scoped conditions; they identify where the deployment claim would
strengthen if the corresponding research-program work landed.

\subsection{Conditional structure of Theorem 1}
\label{sec:conditional_structure}

Theorem~1 binds when its conditions hold.  After the companion
paper integration, the conditions divide into three structurally
distinct categories:

\paragraph{Operational conditions, verifiable from deployment state.}
$I_1, \ldots, I_{11}$, (C1)--(C5), (C5.SPRT) sub-clauses, (C5.MULT),
(C8)--(C12).  Each is computable from the verification ledger or
from auditor-attested deployment-class parameters.

\paragraph{Structural conditions, verifiable from
verification-protocol design.}  (C7) bounded co-evolution (now
derived in \citep[Theorem~1]{lasser2026foundations} from
(C5.HOEFF), (C7.RATE), and channel-projection structure);
Assumption~\ref{ass:concgap_conditions}'s (REP) + (DISP) + (COAL)
+ (EMB) (verifiable from counterparty-selection diversity,
dispersion estimation, coalition-detection audits, and
embedding-$\phi$ certification per
\citep[\S 6]{lasser2026foundations}).

\paragraph{Scope restrictions.}  The five named residuals
(R1)--(R5) bound what the theorem does not claim.  Deployments
falling into any residual are outside the deployment claim's
guarantee; the residuals are explicit, not hidden.

The deployment claim is the conjunction of all three categories.
The companion paper supplies the second category (structural
conditions) as operationally auditable verification-protocol-design
conditions.
